\documentclass[12pt]{article}
\usepackage{../t068report}
% Auto-generated by make_numbers.py from results/*.json — do not edit by hand
\newcommand{\pOneK}{4}
\newcommand{\pOneFS}{pca}
\newcommand{\pOneScaler}{quantile}
\newcommand{\pOneD}{28}
\newcommand{\perComp}{434}
\newcommand{\testRate}{0.130\%}
\newcommand{\testFraud}{74}
\newcommand{\gmmTestAP}{0.686}
\newcommand{\aeTestAP}{0.618}
\newcommand{\daeTestAP}{0.661}
\newcommand{\aeValAP}{0.584}
\newcommand{\daeValAP}{0.680}
\newcommand{\gmmRecallA}{75.7\%}
\newcommand{\gmmPA}{90\%}
\newcommand{\aeRecallA}{70.3\%}
\newcommand{\aePA}{72\%}
\newcommand{\daeRecallA}{70.3\%}
\newcommand{\daePA}{84\%}
\newcommand{\aeSeedMean}{0.569}
\newcommand{\aeSeedStd}{0.068}
\newcommand{\daeSeedMean}{0.628}
\newcommand{\daeSeedStd}{0.013}
\newcommand{\aeScaler}{standard}
\newcommand{\aeLatent}{16}
\newcommand{\aeHidden}{24--16}
\newcommand{\aeAct}{elu}
\newcommand{\aeDropout}{0.0}
\newcommand{\aeWD}{1e-05}
\newcommand{\daeCorruption}{Gaussian noise $\sigma=0.2$}
\newcommand{\nAblations}{27}
\newcommand{\ablDefault}{0.078$\pm$0.021}
\newcommand{\ablationRows}{%
default & default AE & 0.078$\pm$0.021 & 200 \\
A0 input scaling & robust & 0.144$\pm$0.063 & 200 \\
A0 input scaling & quantile & 0.055$\pm$0.004 & 200 \\
A1 bottleneck & m=2 & 0.068$\pm$0.006 & 150 \\
A1 bottleneck & m=4 & 0.101$\pm$0.077 & 161 \\
A1 bottleneck & m=12 & 0.172$\pm$0.083 & 182 \\
A1 bottleneck & m=16 & 0.442$\pm$0.042 & 175 \\
A2 depth & () & 0.034$\pm$0.001 & 45 \\
A2 depth & (16,) & 0.056$\pm$0.002 & 168 \\
A2 depth & (28, 24, 16) & 0.301$\pm$0.184 & 187 \\
A3 noise $\sigma$ & $\sigma$=0.05 & 0.071$\pm$0.020 & 200 \\
A3 noise $\sigma$ & $\sigma$=0.1 & 0.060$\pm$0.008 & 180 \\
A3 noise $\sigma$ & $\sigma$=0.2 & 0.433$\pm$0.037 & 200 \\
A3 noise $\sigma$ & $\sigma$=0.4 & 0.276$\pm$0.045 & 198 \\
A4 masking & p=0.1 & 0.172$\pm$0.099 & 152 \\
A4 masking & p=0.2 & 0.199$\pm$0.002 & 106 \\
A5 dropout & dropout=0.1 & 0.137$\pm$0.005 & 108 \\
A5 dropout & dropout=0.2 & 0.055$\pm$0.003 & 75 \\
A5 weight decay & wd=0.0 & 0.071$\pm$0.011 & 192 \\
A5 weight decay & wd=0.0001 & 0.204$\pm$0.059 & 200 \\
A5 weight decay & wd=0.001 & 0.311$\pm$0.021 & 200 \\
A6 activation & relu & 0.235$\pm$0.024 & 200 \\
A6 activation & tanh & 0.066$\pm$0.003 & 200 \\
A6 activation & linear & 0.034$\pm$0.001 & 18 \\
A7 training data & 10\% of legit rows & 0.083$\pm$0.004 & 200 \\
A7 training data & 50\% of legit rows & 0.198$\pm$0.072 & 200 \\
A7 training data & contaminated (frauds kept) & 0.012$\pm$0.003 & 200 \\
}
\newcommand{\aeMainRows}{%
GMM (Phase 1) & 0.686 & [0.57, 0.81] & 0.934 & 0.757 & 0.784 & 0.784 & 0.90 & 0.56 \\
Undercomplete AE & 0.618 & [0.50, 0.74] & 0.938 & 0.703 & 0.770 & 0.770 & 0.72 & 0.52 \\
Denoising AE & 0.661 & [0.55, 0.78] & 0.940 & 0.703 & 0.784 & 0.797 & 0.84 & 0.55 \\
}
\newcommand{\candidateRows}{%
default & 0.078$\pm$0.021 & ablation study \\
A1 bottleneck: m=16 & 0.442$\pm$0.042 & ablation study \\
A5 weight decay: wd=0.001 & 0.311$\pm$0.021 & ablation study \\
A2 depth: (28, 24, 16) & 0.301$\pm$0.184 & ablation study \\
A6 activation: relu & 0.235$\pm$0.024 & ablation study \\
A0 input scaling: robust & 0.144$\pm$0.063 & ablation study \\
A5 dropout: dropout=0.1 & 0.137$\pm$0.005 & ablation study \\
top-2 combined & 0.275$\pm$0.183 & trained for selection \\
all improvements combined & 0.024$\pm$0.000 & trained for selection \\
chosen AE + $\sigma$=0.2 & 0.621$\pm$0.036 & trained for selection \\
chosen AE + p=0.2 & 0.345$\pm$0.040 & trained for selection \\
}
\newcommand{\aeSigRows}{%
Undercomplete AE $-$ GMM & -0.068 & [-0.107, -0.032] & 1.000 \\
Denoising AE $-$ GMM & -0.025 & [-0.055, +0.001] & 0.964 \\
Denoising AE $-$ Undercomplete AE & +0.043 & [+0.018, +0.075] & 0.000 \\
}
\newcommand{\sigAEGMM}{-0.068 [-0.107, -0.032]}
\newcommand{\sigDAEGMM}{-0.025 [-0.055, +0.001]}
\newcommand{\sigDAEAE}{+0.043 [+0.018, +0.075]}
\newcommand{\diagRhoLegit}{0.429}
\newcommand{\diagRhoFraud}{0.941}
\newcommand{\diagMSE}{0.583}
\newcommand{\diagMaha}{0.171}
\newcommand{\diagReconLegit}{73.6\%}
\newcommand{\diagReconFraud}{31.7\%}
\newcommand{\hybridRows}{%
Phase-1 GMM (all legit rows) & 4 & 0.695$\pm$0.000 & 0.686$\pm$0.000 \\
H0 GMM on x (same rows) & 2 & 0.688$\pm$0.001 & 0.697$\pm$0.013 \\
H3 augmented GMM [x, log e] & 2 & 0.686$\pm$0.002 & 0.697$\pm$0.013 \\
H4 late fusion (rank avg) & -- & 0.342$\pm$0.186 & 0.524$\pm$0.137 \\
AE alone (reconstruction error) & -- & 0.163$\pm$0.111 & 0.325$\pm$0.114 \\
H2 DAGMM-style [z, log e, cos] & 3 & 0.025$\pm$0.002 & 0.044$\pm$0.006 \\
H1 latent GMM [z] & 4 & 0.008$\pm$0.002 & 0.012$\pm$0.004 \\
}
\newcommand{\bestHybrid}{H3 augmented GMM [x, log e]}
\newcommand{\recommended}{GMM (Phase 1)}
\newcommand{\evTopLatent}{4.1\%}
\newcommand{\evTopX}{77.0\%}
\newcommand{\evTopAE}{77.0\%}
\newcommand{\evZlegit}{2.91}
\newcommand{\evZfraud}{3.52}
\newcommand{\evRhoFraud}{0.952}
\newcommand{\evRhoLegit}{0.384}
\newcommand{\semiRows}{%
6 & 0.523$\pm$0.236 & 0.531$\pm$0.252 \\
17 & 0.696$\pm$0.073 & 0.692$\pm$0.072 \\
34 & 0.747$\pm$0.029 & 0.759$\pm$0.021 \\
86 & 0.769$\pm$0.013 & 0.772$\pm$0.009 \\
342 & 0.778$\pm$0.000 & 0.771$\pm$0.000 \\
}
\newcommand{\finalRows}{%
GMM (Phase 1) & 0.686 & [0.57, 0.81] & 0.934 & 0.757 & 0.784 & 0.784 & 0.90 & 0.56 \\
Stage-1 DAE alone (half the rows) & 0.301 & [0.24, 0.40] & 0.943 & 0.189 & 0.716 & 0.770 & 0.48 & 0.40 \\
Hybrid H3 augmented GMM [x, log e] & 0.695 & [0.59, 0.81] & 0.961 & 0.757 & 0.757 & 0.770 & 0.92 & 0.56 \\
Supervised GBM, raw features (uses labels) & 0.778 & [0.69, 0.86] & 0.976 & 0.784 & 0.838 & 0.865 & 0.96 & 0.58 \\
Semi-supervised hybrid H5 (uses labels) & 0.771 & [0.68, 0.86] & 0.978 & 0.757 & 0.811 & 0.865 & 0.94 & 0.56 \\
}
\newcommand{\hybSigRows}{%
H3 augmented GMM [x, log e] $-$ GMM & +0.009 & [-0.009, +0.023] & 0.325 \\
Stage-1 DAE alone $-$ GMM & -0.384 & [-0.474, -0.299] & 1.000 \\
H5 $-$ supervised GBM & -0.007 & [-0.045, +0.033] & 0.647 \\
supervised GBM $-$ GMM & +0.092 & [+0.026, +0.145] & 0.001 \\
}
\newcommand{\evidenceRows}{%
frauds above legit 99th pct — latent GMM on z (H1) & 4.1\% \\
frauds above legit 99th pct — GMM on x (H0) & 77.0\% \\
frauds above legit 99th pct — AE error & 77.0\% \\
mean |z| legit & 2.908 \\
mean |z| fraud & 3.524 \\
Spearman(AE error, GMM score) — test frauds & 0.952 \\
Spearman(AE error, GMM score) — test legit & 0.384 \\
}
\newcommand{\sigHA}{+0.009 [-0.009, +0.023]}
\newcommand{\sigHAname}{Best unsupervised hybrid: H3 augmented GMM [x, log e] $-$ GMM}
\newcommand{\sigHB}{-0.384 [-0.474, -0.299]}
\newcommand{\sigHBname}{AE $-$ GMM}
\newcommand{\sigHC}{-0.007 [-0.045, +0.033]}
\newcommand{\sigHCname}{H5 $-$ supervised GBM}
\newcommand{\sigHD}{+0.092 [+0.026, +0.145]}
\newcommand{\sigHDname}{supervised GBM $-$ GMM}
\newcommand{\recRows}{%
FP budget 0.1\% & 94 & 12.3 & 56/74 & 0.596 & 38 \\
FP budget 0.5\% & 221 & 28.9 & 58/74 & 0.262 & 163 \\
FP budget 1.0\% & 399 & 52.1 & 58/74 & 0.145 & 341 \\
review capacity top-50 & 50 & 6.5 & 45/74 & 0.900 & 5 \\
review capacity top-100 & 100 & 13.1 & 56/74 & 0.560 & 44 \\
review capacity top-200 & 200 & 26.1 & 57/74 & 0.285 & 143 \\
review capacity top-500 & 500 & 65.3 & 58/74 & 0.116 & 442 \\
}


\begin{document}

\maketitlepage{Phase 2: Deep Learning Evaluation}{Autoencoder anomaly detection and hybrid ML + DL density models}

%==================================================================
\section*{Abstract}
\addcontentsline{toc}{section}{Abstract}
In Phase~1 we showed that a Gaussian Mixture Model (GMM), fitted only to legitimate card transactions after a quantile
transform, is a strong label-free fraud ranker (test PR-AUC \gmmTestAP{}). We also showed that it works only after the
inputs are forced to look Gaussian, and that the Expectation--Maximisation (EM) fit is unstable in high dimensions. In
Phase~2 we ask whether \emph{learned non-linear} models can close this gap.

We train undercomplete and denoising autoencoders only on legitimate transactions and score each transaction by its
reconstruction error. We run \nAblations{} ablation variants (two seeds each) covering input scaling, bottleneck size,
depth, corruption type and strength, dropout, weight decay, activation function, amount of training data and label
contamination. The factors interact strongly, so the final networks are chosen among validated candidates rather than
by stacking every individually helpful change. The denoising autoencoder reaches a test PR-AUC of \daeTestAP{}. This is
statistically on par with the GMM (paired bootstrap difference \sigDAEGMM{}), and the model is stable across seeds. The
plain undercomplete autoencoder is significantly worse (\aeTestAP{}).

We then test five ML + DL hybrids, all fitted in two stages without leakage: latent-space GMMs in the style of DAGMM, a
reconstruction-augmented GMM, late score fusion, and semi-supervised stacking. None significantly beats the GMM. Our
measurements show why: the autoencoder and the GMM find the same frauds, and the encoder's latent code discards the
anomaly evidence, which lives in the reconstruction residual. A label-efficiency study shows that the label-free GMM is
worth about 17 labelled frauds, and that a supervised model trained on all labels is significantly better. We
therefore recommend deploying the GMM first and moving to a supervised model once a few dozen frauds have been
confirmed.

\medskip
\textbf{Keywords:} fraud detection, anomaly detection, autoencoder, denoising autoencoder, Gaussian mixture model,
hybrid models, ablation study, label efficiency.

\newpage
\tableofcontents
\newpage

%==================================================================
\section{Introduction}

\subsection{Recap of Phase 1}
Phase~1 framed card-fraud triage as novelty detection. A GMM fitted by EM to \emph{legitimate} transactions ranks every
transaction by its negative log-likelihood $-\log p(\mathbf{x})$, so that analysts review the least ``normal-looking''
transactions first. With a quantile (marginal Gaussianisation) transform and $K=\pOneK$ full-covariance components, the
GMM reached a test PR-AUC of \gmmTestAP{}. It caught \gmmRecallA{} of test frauds at a 0.1\% false-positive budget, and
\gmmPA{} of its top-50 alerts were fraud.

\subsection{The gap left by Phase 1}
Phase~1 also measured three weaknesses of the GMM:
\begin{enumerate}
\item \textbf{Shape assumption.} The GMM works only after every feature is forced to look Gaussian. Each component is a
fixed ellipsoid, and adding components did not help.
\item \textbf{Instability in high dimensions.} Each full-covariance component has \perComp{} parameters in \pOneD{}
dimensions, the fitted covariances are close to singular, and EM's solution depends on its random start.
\item \textbf{Model selection.} AIC and BIC keep asking for more components without any gain in detection quality.
\end{enumerate}

\subsection{Objectives of Phase 2}
\begin{enumerate}
\item Train autoencoders (undercomplete and denoising) on legitimate transactions only, as named in the project brief,
and score transactions by reconstruction error.
\item Justify every architectural and regularisation choice with ablation studies.
\item Compare the autoencoders with the Phase-1 GMM fairly: same rows, same split, same metrics, thresholds from
validation, and paired statistical tests.
\item Design and evaluate hybrid ML + DL approaches, and explain when and why they add value.
\item Give a practical triage recommendation.
\end{enumerate}

\subsection{Contributions}
\begin{itemize}
\item A fully ablated autoencoder study (\nAblations{} variants, two seeds each), with a candidate-based selection
procedure that handles interacting factors.
\item A diagnosis of the relationship between reconstruction error and Gaussian density, tested empirically.
\item A controlled hybrid study with leakage-free two-stage fitting: DAGMM-style latent densities, a
reconstruction-augmented GMM, late fusion, and a semi-supervised stacking model with a label-efficiency analysis.
\item Multi-seed variance and paired bootstrap significance tests for every headline comparison.
\end{itemize}

%==================================================================
\section{Deep Learning Literature Review}\label{sec:lit}

\subsection{Autoencoders and their relation to PCA}
An autoencoder is a neural network trained to reproduce its input through a narrow intermediate layer. Baldi and
Hornik~\cite{baldi1989} proved that a \emph{linear} autoencoder trained with squared error has no spurious local
minima and that its optimum spans the principal subspace of the data. A linear autoencoder is therefore equivalent to
PCA. Hinton and Salakhutdinov~\cite{hinton2006} showed that deep non-linear autoencoders learn much better
low-dimensional codes than PCA. Goodfellow, Bengio and Courville~\cite{goodfellow2016} give a textbook treatment of
undercomplete, regularised and denoising autoencoders.

\subsection{Reconstruction error as an anomaly score}
Sakurada and Yairi~\cite{sakurada2014} train autoencoders on normal data and use reconstruction error as the anomaly
score. They show that this captures non-linear structure that linear PCA misses. An and Cho~\cite{an2015} propose a
variational autoencoder with a ``reconstruction probability'' score. The surveys of Chalapathy and
Chawla~\cite{chalapathy2019} and Pang~\emph{et~al.}~\cite{pang2021} identify reconstruction-based models as the dominant
family of deep anomaly detectors. They also point out their main weakness: a network that generalises well can
reconstruct anomalies too. Pumsirirat and Yan~\cite{pumsirirat2018} apply autoencoders to credit-card fraud.

\subsection{Denoising and regularisation}
Vincent~\emph{et~al.}~\cite{vincent2008,vincent2010} introduced denoising autoencoders, which reconstruct a clean
input from a corrupted copy. This forces robust features and prevents the network from learning the identity map.
Alain and Bengio~\cite{alain2014} showed that, for small Gaussian noise, the optimal denoising reconstruction satisfies
$r(\mathbf{x})-\mathbf{x}\propto\sigma^2\nabla_{\mathbf{x}}\log p(\mathbf{x})$. A denoising autoencoder therefore
implicitly learns the gradient of the log-density of the data, which links the reconstruction view of Phase~2 to the
density view of Phase~1. Our training also uses dropout~\cite{srivastava2014}, the Adam optimiser~\cite{kingma2015} and
ELU activations~\cite{clevert2016}, each of which we ablate.

\subsection{Hybrid deep and density models}
DAGMM~\cite{zong2018} couples an autoencoder with a GMM fitted on the concatenation of the latent code and two
reconstruction features (relative Euclidean distance and cosine similarity). It trains both jointly and outperforms
stand-alone autoencoders and GMMs on several benchmarks. Deep SVDD~\cite{ruff2018} learns a representation in which
normal data falls inside a minimum-volume hypersphere. For card fraud, Carcillo~\emph{et~al.}~\cite{carcillo2021} combine
unsupervised outlier scores with supervised classifiers.

\subsection{Summary}
\begin{table}[H]
\centering\small
\caption{Deep-learning literature and how it shapes our Phase-2 design.}
\label{tab:lit}
\begin{tabularx}{\linewidth}{@{}l X X@{}}
\toprule
\textbf{Work} & \textbf{Contribution} & \textbf{How we use it}\\
\midrule
Baldi \& Hornik~\cite{baldi1989} & Linear autoencoder = PCA & Ablation of the activation function (linear case)\\
Sakurada \& Yairi~\cite{sakurada2014} & Reconstruction error as anomaly score & Phase-2 baseline detector\\
Vincent \emph{et al.}~\cite{vincent2008} & Denoising autoencoder & Denoising variant; noise ablation\\
Alain \& Bengio~\cite{alain2014} & Denoising learns the density gradient & Link between the GMM and the autoencoder\\
Zong \emph{et al.}~\cite{zong2018} (DAGMM) & GMM on latent code + reconstruction features & Hybrids H1 and H2\\
Carcillo \emph{et al.}~\cite{carcillo2021} & Unsupervised scores as features for supervised models & Semi-supervised hybrid H5\\
Pang \emph{et al.}~\cite{pang2021} & Reconstruction models can over-generalise & Diagnosis of over-generalisation\\
\bottomrule
\end{tabularx}
\end{table}

%==================================================================
\section{Data, Protocol and DL-specific Data Handling}\label{sec:data}

\subsection{Shared protocol}
All Phase-2 models use exactly the same data and protocol as Phase~1:
\begin{itemize}
\item 1{,}081 duplicate rows are removed before splitting;
\item rows are ordered by time and split into train (60\%), validation (20\%) and test (20\%) blocks;
\item models are fitted on legitimate training rows only;
\item hyper-parameters and thresholds are chosen on validation;
\item the test block (\testFraud{} frauds, base rate \testRate) is used once.
\end{itemize}
The input features are the Phase-1 feature set (V1--V28, $d=\pOneD$).

\subsection{Data handling for the deep models}
\begin{itemize}
\item \textbf{Legitimate-only training.} The fraud label is used once, to remove the 342 training frauds. We also
measure what happens when they are left in (ablation A7).
\item \textbf{Label-free early stopping.} A random 10\% of the legitimate training rows is held out and used only to
monitor the reconstruction loss. No fraud label is involved in training or early stopping.
\item \textbf{Input scaling.} Each model gets its own validated preprocessing. The GMM uses the quantile transform from
Phase~1. For the autoencoder the scaler is chosen by ablation (A0); every scaler is fitted on legitimate training rows
only.
\item \textbf{Leakage control for hybrids.} Second-stage models are never fitted on rows that the autoencoder was
trained on (Section~\ref{sec:hyb}).
\end{itemize}

%==================================================================
\section{Methodology}\label{sec:method}

\subsection{Autoencoder}
An encoder $f_\theta:\mathbb{R}^d\rightarrow\mathbb{R}^m$ and a decoder $g_\phi:\mathbb{R}^m\rightarrow\mathbb{R}^d$ with
$m<d$ are trained on legitimate rows to minimise the mean squared reconstruction error plus an L2 penalty:
\begin{equation}
\mathcal{L}(\theta,\phi)=\frac1n\sum_{i=1}^{n}\big\lVert\mathbf{x}_i-g_\phi\big(f_\theta(\tilde{\mathbf{x}}_i)\big)\big\rVert_2^2+\lambda\lVert\mathbf{W}\rVert_2^2 .
\end{equation}
For the undercomplete autoencoder $\tilde{\mathbf{x}}_i=\mathbf{x}_i$. The bottleneck ($m<d$) stops the network from
learning the identity map: it must keep only the structure shared by legitimate transactions, the (possibly curved)
manifold on which they lie.

\subsection{Denoising autoencoder}
The denoising autoencoder (DAE) trains on corrupted inputs and reconstructs the \emph{clean} input. We consider two
corruptions:
\begin{align}
\text{Gaussian noise:}&\quad \tilde{\mathbf{x}}=\mathbf{x}+\bm\epsilon,\qquad \bm\epsilon\sim\mathcal{N}(\mathbf{0},\sigma^2\mathbf{I}),\\
\text{Masking:}&\quad \tilde{x}_j=x_j\, b_j,\qquad b_j\sim\text{Bernoulli}(1-p).
\end{align}
The network must learn to pull points back onto the data manifold. This acts as a strong regulariser, and by the
result of Alain and Bengio~\cite{alain2014} it ties the reconstruction to the gradient of the data's log-density. At
test time no noise is added.

\subsection{Anomaly score}
Every transaction is scored by its reconstruction error on the clean input,
\begin{equation}
e(\mathbf{x})=\frac1d\lVert\mathbf{x}-\hat{\mathbf{x}}\rVert_2^2,\qquad \hat{\mathbf{x}}=g_\phi(f_\theta(\mathbf{x})).
\end{equation}
Thresholds are chosen on validation exactly as in Phase~1.

\subsection{Architecture}
The default network is a symmetric multilayer perceptron (Fig.~\ref{fig:ae}). Table~\ref{tab:arch} lists the design
choices and their reasons.

\begin{figure}[H]
\centering
\begin{tikzpicture}[x=1.35cm,y=1cm]
\foreach \i/\h/\lab in {0/3.0/28,1/2.6/24,2/1.9/16,3/1.9/16,4/1.9/16,5/2.6/24,6/3.0/28}{
  \draw (\i-0.22,-\h/2) rectangle (\i+0.22,\h/2);
  \node[below,font=\small] at (\i,-1.7) {\lab};
}
\foreach \i in {0,...,5}{\draw[-{Stealth[length=2mm]}] (\i+0.27,0) -- (\i+0.73,0);}
\node[font=\small] at (0,1.9) {input $\mathbf{x}$};
\node[font=\small] at (3,1.4) {code $\mathbf{z}$};
\node[font=\small] at (6,1.9) {output $\hat{\mathbf{x}}$};
\node[font=\small] at (1,-2.3) {encoder};
\node[font=\small] at (5,-2.3) {decoder};
\end{tikzpicture}
\caption{Selected autoencoder architecture ($d=\pOneD$ inputs, hidden layers of 24 and 16 units, a 16-unit code).
Numbers give the width of each layer. The denoising version uses the same network with Gaussian input noise during
training.}
\label{fig:ae}
\end{figure}

\begin{table}[H]
\centering\small
\caption{Architecture and training choices.}
\label{tab:arch}
\begin{tabularx}{\linewidth}{@{}l >{\raggedright\arraybackslash}X >{\raggedright\arraybackslash}X@{}}
\toprule
\textbf{Choice} & \textbf{Value} & \textbf{Reason}\\
\midrule
Encoder & $d\rightarrow24\rightarrow16\rightarrow m$ & gradual compression; $m$ chosen by ablation (A1)\\
Decoder & mirror of the encoder, linear output & standardised inputs are unbounded, so no sigmoid or tanh at the output\\
Code layer & linear & lets the code use the whole real line\\
Activation & ELU & smooth and non-zero for negative inputs (no ``dead'' units); compared in A6\\
Loss & mean squared error & matches the reconstruction score\\
Regularisation & bottleneck, input noise, L2 weight decay, dropout, early stopping & each ablated (A1, A3--A5)\\
Optimiser & Adam, learning rate $2\times10^{-3}$, batch 2048 & standard and stable for small networks\\
Schedule & ReduceLROnPlateau (factor 0.5, patience 3), at most 200 epochs & lowers the learning rate when the holdout loss stalls\\
Early stopping & patience 10 on a label-free 10\% holdout & best weights are restored\\
\bottomrule
\end{tabularx}
\end{table}

\begin{algorithm}[H]
\caption{Training an autoencoder on legitimate transactions}
\label{alg:train}
\begin{algorithmic}[1]
\Require legitimate training rows $X$; configuration (architecture, noise $\sigma$ or masking $p$, weight decay $\lambda$)
\State Split $X$ at random into a training part (90\%) and a label-free holdout $H$ (10\%)
\For{epoch $=1,\dots,200$}
  \For{each mini-batch $B$ of the training part}
    \State $\tilde B\gets$ corrupt$(B)$ \Comment{identity for the undercomplete autoencoder}
    \State $\hat B\gets g_\phi(f_\theta(\tilde B))$;\; loss $\gets\lVert \hat B-B\rVert^2/|B|$
    \State update $\theta,\phi$ with Adam (L2 weight decay $\lambda$)
  \EndFor
  \State holdout loss $\gets$ mean squared error of the clean reconstruction of $H$
  \State keep the weights with the lowest holdout loss; stop after 10 epochs without improvement
\EndFor
\State \Return best weights; score $e(\mathbf{x})=\lVert\mathbf{x}-\hat{\mathbf{x}}\rVert^2/d$
\end{algorithmic}
\end{algorithm}

\subsection{Why reconstruction error may differ from a Gaussian density}
Let the covariance of legitimate data have eigenpairs $(\lambda_j,\mathbf{u}_j)$. The Gaussian score (Mahalanobis
distance) is $\sum_j\big(\mathbf{u}_j^{\top}(\mathbf{x}-\bm\mu)\big)^2/\lambda_j$. It strongly amplifies deviations along
directions in which legitimate data barely varies (small $\lambda_j$). A linear autoencoder with code size $m$ has
error $\sum_{j>m}\big(\mathbf{u}_j^{\top}(\mathbf{x}-\bm\mu)\big)^2$: the residual directions are weighted
\emph{equally}, regardless of how tightly legitimate data is constrained along them. A non-linear autoencoder inherits
this equal weighting in its residual. It may also partly reconstruct anomalies (over-generalisation~\cite{pang2021}).
We test both hypotheses in Section~\ref{sec:diag}.

\subsection{Hybrid designs}\label{sec:hyb}
We evaluate five ways of combining the GMM (machine learning) with the autoencoder (deep learning). Each tests one idea
(Table~\ref{tab:hyb-design}, Fig.~\ref{fig:hyb-pipe}).

\begin{table}[H]
\centering\small
\caption{Hybrid designs.}
\label{tab:hyb-design}
\begin{tabularx}{\linewidth}{@{}l l X@{}}
\toprule
\textbf{ID} & \textbf{Stage-2 model} & \textbf{Idea}\\
\midrule
H0 & GMM on $\mathbf{x}$ (same rows) & reference: the plain GMM fitted on the same rows as the hybrids\\
H1 & GMM on $\mathbf{z}$ & density estimation may be easier in a low-dimensional learned space (DAGMM idea~\cite{zong2018})\\
H2 & GMM on $[\mathbf{z},\log e,\cos(\mathbf{x},\hat{\mathbf{x}})]$ & combine \emph{where} a point is on the manifold with \emph{how far} it is from it (DAGMM)\\
H3 & GMM on $[\mathbf{x},\log e]$ & keep all input information and add the learned error as an extra dimension\\
H4 & rank average of GMM and AE scores & naive late fusion, the baseline any real hybrid must beat\\
H5 & gradient boosting on $[\mathbf{x},\text{GMM score},\log e]$ & semi-supervised: use the few frauds analysts have labelled~\cite{carcillo2021}\\
\bottomrule
\end{tabularx}
\end{table}

\begin{figure}[H]
\centering
\begin{tikzpicture}[node distance=0.8cm and 0.55cm,
  box/.style={draw,align=center,font=\small,minimum height=1.1cm,text width=2.45cm},
  arr/.style={-{Stealth[length=2mm]}}]
\node[box] (a) {Legitimate rows,\\half A};
\node[box,right=of a] (ae) {Stage 1:\\autoencoder};
\node[box,below=1.2cm of a] (b) {Legitimate rows,\\half B (40{,}000)};
\node[box,right=of b] (feat) {Stage-2 features\\$\mathbf{z}$, $\log e$, $\cos$, $\mathbf{x}$};
\node[box,right=of feat] (gmm) {Stage 2:\\GMM ($K$ on val.)};
\node[box,right=of gmm] (s) {Score\\$-\log p(\mathbf{h})$};
\draw[arr] (a) -- (ae);
\draw[arr] (ae) -- (feat);
\draw[arr] (b) -- (feat);
\draw[arr] (feat) -- (gmm);
\draw[arr] (gmm) -- (s);
\end{tikzpicture}
\caption{Two-stage fitting of the unsupervised hybrids (H1--H3). The autoencoder is trained on half A; every
second-stage GMM is fitted on half B, rows the autoencoder never saw.}
\label{fig:hyb-pipe}
\end{figure}

\textbf{Leakage control.} On the rows it was trained on, an autoencoder's reconstruction error is unrealistically
small. A second-stage model fitted on those rows would learn a distorted ``normal'' error level. For H1--H4 the
legitimate training rows are therefore split at random into two halves. Half A trains the autoencoder, and a random
40{,}000 rows of half B fit every second-stage GMM, together with the H0 reference on the same rows. The whole procedure
is repeated for three seeds, and for each second-stage GMM, $K\in\{1,2,3,4\}$ is chosen on validation. For H5, GMM and
autoencoder scores of training rows are computed \emph{out of fold}: in a 2-fold cross-fitting, each half is scored by
models fitted on the other half. The supervised model is trained with a varying number of labelled frauds (5 random
draws per budget).

\textbf{Why two-stage instead of joint training?} DAGMM trains both parts jointly, which is known to be fragile: the GMM
energy term can pull the encoder into degenerate solutions. Two-stage fitting is stable, and each stage can be
validated separately. Joint training is left for future work.

\subsection{Evaluation and statistical validation}
We use the Phase-1 metrics:
\begin{itemize}
\item PR-AUC, with a 95\% stratified bootstrap interval;
\item recall at false-positive budgets of 0.1\%, 0.5\% and 1\% (thresholds from validation);
\item precision at review capacity (P@$k$).
\end{itemize}
To decide whether two models really differ, we use a \emph{paired} stratified bootstrap. We resample the test set
1{,}000 times, compute $\text{PR-AUC}(A)-\text{PR-AUC}(B)$ on each identical resample, and report the 95\% interval of the
difference together with the share of resamples in which $A$ is not better than $B$. Every development decision uses
validation PR-AUC averaged over seeds.

\subsection{Experimental setup}
\begin{table}[H]
\centering\small
\caption{Experimental setup.}
\label{tab:setup}
\begin{tabularx}{\linewidth}{@{}l X@{}}
\toprule
\textbf{Item} & \textbf{Setting}\\
\midrule
Ablations & \nAblations{} variants $\times$ 2 seeds; one factor changed at a time from the default network\\
Final models & 3 seeds each; test scores averaged over seeds\\
Hybrids & 3 seeds; second-stage $K\in\{1,2,3,4\}$ chosen on validation\\
Label efficiency & 6, 17, 34, 86 and 342 labelled frauds; 5 random draws per budget\\
Gradient boosting & 300 iterations, learning rate 0.05, 15 leaves, L2 regularisation 1.0\\
Bootstrap & 1{,}000 stratified resamples (paired for model comparisons)\\
Software & Python 3, PyTorch~\cite{paszke2019}, scikit-learn~\cite{pedregosa2011}\\
Hardware & Google Colab with an NVIDIA T4 GPU; about 35 s per autoencoder training run\\
\bottomrule
\end{tabularx}
\end{table}

%==================================================================
\section{Experiments and Results}\label{sec:results}

\subsection{Ablation studies}\label{sec:abl}
Fig.~\ref{fig:abl} shows all ablation groups, and Table~\ref{tab:abl} (Appendix~\ref{app:abl}) lists every variant.
The default network scores only \ablDefault{} on validation. The main findings are:
\begin{enumerate}
\item \textbf{Bottleneck size is the largest single effect (A1).} Codes of 2--8 units are poor, and $m=16$ is by far
the best variant ($\approx0.44$). The V-features are principal components with fairly evenly spread variance, so
legitimate behaviour has no very low-dimensional structure. A tight code cannot rebuild \emph{legitimate} rows well,
and their large errors drown the fraud signal.
\item \textbf{Non-linearity is essential (A6).} The linear autoencoder, which is equivalent to PCA~\cite{baldi1989}, is
the worst variant ($\approx0.03$). A network with no hidden layer also fails (A2).
\item \textbf{Denoising helps when the noise is right (A3, A4).} Gaussian noise of $\sigma=0.2$ gives a large gain;
smaller noise does nothing and $\sigma=0.4$ is too strong. Masking noise is weaker than Gaussian noise here.
\item \textbf{Regularisation (A5).} Dropout \emph{hurts}, because it makes reconstructions noisier, and that noise is
exactly what we score. Stronger L2 weight decay ($10^{-3}$) helps by giving a smoother mapping that generalises less to
unusual inputs.
\item \textbf{Inputs (A0).} Unlike the GMM, the autoencoder does not benefit from quantile inputs, which clip the extreme
values that expose fraud. Robust scaling is slightly better than standard scaling on its own.
\item \textbf{Training data (A7).} More legitimate data helps. Leaving the 0.2\% of training frauds in the training set
is catastrophic, because the network learns to rebuild those extreme points too. The legitimate-only protocol is
essential for this model.
\end{enumerate}

\begin{figure}[H]
\centering
\includegraphics[width=\linewidth]{03_ablations.png}
\caption{Autoencoder ablations: validation PR-AUC (mean $\pm$ std over 2 seeds) per factor. The solid line is the
default network; the dashed line is the Phase-1 GMM.}
\label{fig:abl}
\end{figure}

\subsection{Model selection with interacting factors}
A natural next step would be to combine every change that helped on its own. We tried this: the combined network
reached a validation PR-AUC of only about 0.02, worse than the default. The factors \emph{interact}. For example, a
wide code combined with dropout and a deeper stack behaves very differently from each change alone. We therefore
compare a short list of candidates on validation (Table~\ref{tab:cand}):
\begin{itemize}
\item the default network;
\item each factor's best setting on its own;
\item the two strongest changes combined;
\item all improving changes combined.
\end{itemize}
The winner is the single best change, a 16-unit code with \aeScaler{} inputs. For the denoising autoencoder we add the
best Gaussian-noise and the best masking setting to this network and again compare on validation. Gaussian noise
$\sigma=0.2$ wins clearly.

\begin{table}[H]
\centering\small
\caption{Candidate configurations compared on validation (PR-AUC, mean $\pm$ std over 2 seeds).}
\label{tab:cand}
\begin{tabular}{@{}l c l@{}}
\toprule
\textbf{Candidate} & \textbf{Validation PR-AUC} & \textbf{Source}\\
\midrule
\candidateRows
\bottomrule
\end{tabular}
\end{table}

\subsection{Training behaviour}
Fig.~\ref{fig:curves} shows the training and holdout losses of the final networks. Both converge smoothly; early
stopping on the label-free holdout restores the best epoch. For the denoising autoencoder the training loss is measured
on corrupted inputs, so it stays above the clean holdout loss.

\begin{figure}[H]
\centering
\includegraphics[width=0.95\linewidth]{03_training_curves.png}
\caption{Training loss and label-free holdout loss of the final autoencoder (left) and denoising autoencoder (right),
seed 0.}
\label{fig:curves}
\end{figure}

\subsection{Autoencoders versus the Phase-1 GMM}
Table~\ref{tab:ae} compares the final models on the test set (autoencoder scores averaged over three seeds),
Table~\ref{tab:aesig} gives the paired significance tests, and Fig.~\ref{fig:aepr} shows the curves.
\begin{itemize}
\item The \textbf{denoising autoencoder} reaches PR-AUC \daeTestAP{}. It is statistically on par with the GMM
(\gmmTestAP{}; difference \sigDAEGMM{}) and stable across seeds (\daeSeedMean{} $\pm$ \daeSeedStd{}).
\item The \textbf{undercomplete autoencoder} is significantly worse than the GMM (\sigAEGMM{}) and sensitive to the seed
(\aeSeedMean{} $\pm$ \aeSeedStd{}).
\item \textbf{Denoising improves on plain reconstruction} by \sigDAEAE{}, acting as a strong
regulariser~\cite{vincent2008}.
\item At strict budgets the GMM remains ahead. It catches \gmmRecallA{} of frauds at a 0.1\% false-positive rate against
\daeRecallA{} for the denoising autoencoder, and its top-50 queue holds \gmmPA{} fraud against \daePA{}.
\end{itemize}

\begin{table}[H]
\centering\small
\caption{Test results: Phase-1 GMM versus autoencoders (thresholds from validation; base rate \testRate).}
\label{tab:ae}
\setlength{\tabcolsep}{3pt}
\resizebox{\linewidth}{!}{%
\begin{tabular}{@{}l c c c c c c c c@{}}
\toprule
\textbf{Model} & \textbf{PR-AUC} & \textbf{95\% CI} & \textbf{ROC-AUC} & \textbf{Rec@0.1\%} & \textbf{Rec@0.5\%} & \textbf{Rec@1\%} & \textbf{P@50} & \textbf{P@100}\\
\midrule
\aeMainRows
\bottomrule
\end{tabular}}
\end{table}

\begin{table}[H]
\centering\small
\caption{Paired bootstrap tests on the test set (1{,}000 resamples).}
\label{tab:aesig}
\begin{tabular}{@{}l r c c@{}}
\toprule
\textbf{Comparison} & \textbf{$\Delta$ PR-AUC} & \textbf{95\% CI} & \textbf{P(A not better)}\\
\midrule
\aeSigRows
\bottomrule
\end{tabular}
\end{table}

\begin{figure}[H]
\centering
\includegraphics[width=\linewidth]{03_pr_curves_gmm_vs_ae.png}
\caption{Test precision--recall curves (left) and recall versus false-positive rate (right) for the GMM and the two
autoencoders.}
\label{fig:aepr}
\end{figure}

\subsection{Threshold analysis}
Fig.~\ref{fig:aetrade} shows precision and recall against review workload for the three models. The curves almost
coincide for large queues. The GMM is slightly better for small, high-precision queues, which matter most when
analysts have little capacity.

\begin{figure}[H]
\centering
\includegraphics[width=0.85\linewidth]{03_threshold_tradeoff.png}
\caption{Precision (dashed) and recall (solid) versus number of alerts on the test period.}
\label{fig:aetrade}
\end{figure}

\subsection{Interpretability: which features does the autoencoder fail to rebuild?}
Reconstruction error decomposes over features, $e(\mathbf{x})=\frac1d\sum_j(x_j-\hat x_j)^2$. Averaging each term over
frauds and over legitimate rows shows which features the network cannot rebuild for fraud (Fig.~\ref{fig:feat}). V17,
V14 and V12 dominate, the same components that separate the classes in the exploratory analysis of Phase~1. An analyst
can therefore be given a reason for each alert. Fig.~\ref{fig:latent} shows the score distribution and a 2-D
projection of the latent codes.

\begin{figure}[H]
\centering
\includegraphics[width=0.85\linewidth]{03_feature_errors.png}
\caption{Mean squared reconstruction error per feature for legitimate and fraudulent test transactions (log scale).}
\label{fig:feat}
\end{figure}

\begin{figure}[H]
\centering
\includegraphics[width=\linewidth]{03_scores_and_latent.png}
\caption{Autoencoder score distribution on test (left) and the latent codes projected to two dimensions (right).}
\label{fig:latent}
\end{figure}

\subsection{Diagnosis}\label{sec:diag}
Three measurements test the hypotheses of Section~\ref{sec:method}.
\begin{enumerate}
\item \textbf{Agreement.} On frauds, the autoencoder error and the GMM score rank transactions almost identically
(Spearman $\rho=\diagRhoFraud$). On legitimate rows they agree much less ($\rho=\diagRhoLegit$). The two models find the
same frauds but raise \emph{different} false alarms. A hybrid could therefore only hope to reduce false alarms, not to
find new frauds.
\item \textbf{Direction weighting.} Re-scoring the same autoencoder's residuals with a Gaussian fitted to legitimate
residuals (a Mahalanobis residual score) \emph{lowers} test PR-AUC from \diagMSE{} to \diagMaha{}. The residuals are
themselves heavy-tailed, so the theoretically motivated fix fails for the same reason a GMM fails on raw heavy-tailed
inputs. The hypothesis is not supported in this simple form.
\item \textbf{Over-generalisation.} The autoencoder reproduces \diagReconLegit{} of a typical legitimate row's deviation
from the mean, but only \diagReconFraud{} of a typical fraud's. It does partly rebuild fraud, which limits its recall,
but the gap is large enough to rank most frauds high.
\end{enumerate}

\subsection{Hybrid study}
Table~\ref{tab:hyb} reports each unsupervised hybrid at its best second-stage $K$ (mean $\pm$ std over three seeds),
and Fig.~\ref{fig:hyb} shows the comparison and the effect of $K$.
\begin{itemize}
\item \textbf{H3} (GMM on raw features plus the log reconstruction error) ties the plain GMM fitted on the same rows
(H0). Adding the learned error gives the mixture nothing it does not already know.
\item \textbf{H4} (late fusion) is worse than the GMM alone. The stage-1 network, trained on only half the legitimate
rows, is weak and unstable on its own, and averaging it in dilutes the GMM's ranking.
\item \textbf{H1 and H2}, the DAGMM-style latent-space hybrids, fail almost completely.
\end{itemize}

\begin{table}[H]
\centering\small
\caption{Unsupervised hybrids at the best second-stage $K$ on validation (mean $\pm$ std over 3 seeds).}
\label{tab:hyb}
\begin{tabular}{@{}l c c c@{}}
\toprule
\textbf{Model} & \textbf{$K$} & \textbf{Validation PR-AUC} & \textbf{Test PR-AUC}\\
\midrule
\hybridRows
\bottomrule
\end{tabular}
\end{table}

\begin{figure}[H]
\centering
\includegraphics[width=\linewidth]{04_hybrid_ablation.png}
\caption{Hybrid ablation. Left: validation (bars, mean $\pm$ std) and test (dots) PR-AUC per design. Right: effect of
the second-stage $K$.}
\label{fig:hyb}
\end{figure}

\subsection{Where does the anomaly information live?}
Table~\ref{tab:evidence} explains the failure of the latent-space hybrids. At a 1\% false-positive rate the latent GMM
flags only \evTopLatent{} of test frauds, while the GMM on the raw inputs flags \evTopX{} and the reconstruction error
\evTopAE{}. Fraud codes are barely larger than legitimate ones (mean $|\mathbf{z}|$ \evZfraud{} against \evZlegit{}).
An encoder trained only on legitimate data maps \emph{any} input close to the region of normal codes. What exposes a
fraud is that the decoder cannot rebuild it from that normal-looking code. The evidence therefore lies in the
reconstruction residual, not in the code, and a density model of the code cannot see it.

\begin{table}[H]
\centering\small
\caption{Evidence on the test set (seed 0). ``Above legit 99th pct'' = share of frauds flagged at a 1\% false-positive
rate.}
\label{tab:evidence}
\begin{tabular}{@{}l r@{}}
\toprule
\textbf{Measurement} & \textbf{Value}\\
\midrule
\evidenceRows
\bottomrule
\end{tabular}
\end{table}

\subsection{Semi-supervised hybrid and label efficiency}
Table~\ref{tab:semi} and Fig.~\ref{fig:semi} train gradient boosting on the training block with only $n$ labelled
frauds (five random draws per budget):
\begin{itemize}
\item with 6 labelled frauds the supervised model is worse than the label-free GMM and very unstable;
\item with about 17 it matches the GMM;
\item with 34 or more it is clearly better.
\end{itemize}
In other words, \emph{the label-free GMM is worth roughly 17 labelled frauds}. Adding the GMM score and the
reconstruction error as features (H5) does not change PR-AUC measurably at any budget. Gradient boosting already extracts
the relevant structure from the raw features.

\begin{table}[H]
\centering\small
\caption{Label efficiency: test PR-AUC (mean $\pm$ std over 5 random draws) against the number of labelled training
frauds.}
\label{tab:semi}
\begin{tabular}{@{}r c c@{}}
\toprule
\textbf{Labelled frauds} & \textbf{Raw features} & \textbf{Raw + GMM score + AE error (H5)}\\
\midrule
\semiRows
\bottomrule
\end{tabular}
\end{table}

\begin{figure}[H]
\centering
\includegraphics[width=0.8\linewidth]{04_label_efficiency.png}
\caption{Label efficiency. The dashed line is the label-free GMM.}
\label{fig:semi}
\end{figure}

\subsection{Final comparison}
Table~\ref{tab:final} and Fig.~\ref{fig:final} compare all approaches on the test set, and Table~\ref{tab:hybsig} gives
the paired tests. No unsupervised hybrid beats the Phase-1 GMM: the best one (H3) differs by \sigHA{}, which is not
significant. Labels are what move the needle: supervised gradient boosting with all training labels beats the GMM by
\sigHD{}. The semi-supervised hybrid is no better than supervised alone (\sigHC{}). Note that the ``Stage-1 DAE alone''
row is the autoencoder trained on half the rows for the hybrid study; the full-data denoising autoencoder of
Table~\ref{tab:ae} scores \daeTestAP{}.

\begin{table}[H]
\centering\small
\caption{Final test comparison (thresholds from validation; base rate \testRate). Models marked ``uses labels'' are
references, not fair competitors to the label-free models.}
\label{tab:final}
\setlength{\tabcolsep}{3pt}
\resizebox{\linewidth}{!}{%
\begin{tabular}{@{}l c c c c c c c c@{}}
\toprule
\textbf{Model} & \textbf{PR-AUC} & \textbf{95\% CI} & \textbf{ROC-AUC} & \textbf{Rec@0.1\%} & \textbf{Rec@0.5\%} & \textbf{Rec@1\%} & \textbf{P@50} & \textbf{P@100}\\
\midrule
\finalRows
\bottomrule
\end{tabular}}
\end{table}

\begin{table}[H]
\centering\small
\caption{Paired bootstrap tests for the hybrid study (1{,}000 resamples).}
\label{tab:hybsig}
\begin{tabular}{@{}l r c c@{}}
\toprule
\textbf{Comparison} & \textbf{$\Delta$ PR-AUC} & \textbf{95\% CI} & \textbf{P(A not better)}\\
\midrule
\hybSigRows
\bottomrule
\end{tabular}
\end{table}

\begin{figure}[H]
\centering
\includegraphics[width=\linewidth]{04_final_comparison.png}
\caption{Final comparison on test: precision--recall curves (left) and recall at fixed false-positive budgets (right).}
\label{fig:final}
\end{figure}

\subsection{Triage recommendation}
The recommended label-free system, chosen on validation, is the Phase-1 GMM. Table~\ref{tab:rec} lists its operating
points on the 7.7-hour test period. A strict 0.1\% budget produces about 12 alerts per hour, of which roughly 60\% are
fraud, and catches three quarters of all frauds. A larger queue adds very few frauds at a high cost in false alarms.

\begin{table}[H]
\centering\small
\caption{Operating points of the recommended label-free model on the test period.}
\label{tab:rec}
\begin{tabular}{@{}l r r r r r@{}}
\toprule
\textbf{Policy} & \textbf{Alerts} & \textbf{Per hour} & \textbf{Caught} & \textbf{Precision} & \textbf{Legit flagged}\\
\midrule
\recRows
\bottomrule
\end{tabular}
\end{table}

%==================================================================
\section{Discussion}\label{sec:disc}

\subsection{Technical validation}
Every claim in this report rests on validation-based decisions and a single use of the test set. Headline comparisons
are averaged over several seeds and checked with paired bootstrap tests. All second-stage models are fitted without
leakage. The same pipeline run twice on Colab produced identical results.

\subsection{Why the hybrids did not beat the GMM}
On this dataset the quantile-transformed Gaussian density already captures what the autoencoder learns about fraud.
Both models flag the same frauds, and they differ only on legitimate rows. Hybrids that keep the raw inputs therefore
tie the GMM, and hybrids that rely on the latent code lose the very information that identifies fraud. This is a useful
negative result: it shows \emph{where} the anomaly signal lives (in the residual, not the code), and it shows that
DAGMM-style designs need care when the encoder is trained only on normal data.

\subsection{Practical implications}
For a bank starting a fraud-triage system, the results suggest a simple path:
\begin{enumerate}
\item deploy the label-free GMM first, with a threshold set by the review capacity;
\item collect analyst verdicts on the alerts;
\item once a few dozen frauds are confirmed, switch to (or add) a supervised model;
\item retrain both regularly, because fraud patterns drift.
\end{enumerate}

\subsection{Limitations}
\begin{itemize}
\item The data covers two days, with 74 test frauds, so confidence intervals are wide (about $\pm0.1$ PR-AUC).
\item The features are anonymised principal components, which limits feature engineering and interpretation.
\item The hybrids are trained in two stages rather than jointly.
\item The stage-1 autoencoder of the hybrid study is trained on half the legitimate rows, which weakens late fusion.
\end{itemize}

%==================================================================
\section{Conclusion and Future Work}\label{sec:concl}
Autoencoders trained only on legitimate transactions can detect card fraud, but their design must be validated
carefully. The bottleneck size, the non-linearity and Gaussian input noise decide performance, and stacking
individually helpful changes without re-validation fails. The selected denoising autoencoder matches the Phase-1 GMM
within statistical noise and is the most stable deep model. No hybrid significantly improves on the GMM. The best one
ties it, and DAGMM-style latent densities fail because the encoder maps fraud to normal-looking codes. Labels matter more
than model combination: the label-free GMM is worth about 17 labelled frauds, and supervised boosting with all labels
is significantly better.

Future work includes:
\begin{itemize}
\item end-to-end (joint) DAGMM training;
\item flexible density models, such as normalising flows, on the reconstruction residual rather than the latent code;
\item longer data with real concept drift;
\item cost-sensitive thresholds that use transaction amounts.
\end{itemize}

%==================================================================
\bibliographystyle{IEEEtran}
\bibliography{../references}

%==================================================================
\appendix
\section{Full Ablation Results}\label{app:abl}
\begin{table}[H]
\centering\small
\caption{All autoencoder ablation variants (validation PR-AUC, mean $\pm$ std over 2 seeds; mean number of epochs
before early stopping).}
\label{tab:abl}
\begin{tabular}{@{}l l c r@{}}
\toprule
\textbf{Group} & \textbf{Setting} & \textbf{Validation PR-AUC} & \textbf{Epochs}\\
\midrule
\ablationRows
\bottomrule
\end{tabular}
\end{table}

\section{Reproducibility}
\begin{itemize}
\item \textbf{Code.} Phase 2 uses notebooks 03 (autoencoders) and 04 (hybrid study). Both run top-to-bottom on Google
Colab with a T4 GPU, download the dataset and verify its checksum.
\item \textbf{Seeds.} Fixed seed lists for ablations (2 seeds), final models (3 seeds), hybrids (3 seeds) and label
draws (5 draws).
\item \textbf{Traceability.} All numbers in this report are generated from \texttt{results/*.json} into
\texttt{numbers.tex}; all figures come from \texttt{figures/}.
\item \textbf{Run time on Colab.} Notebook 03 takes about 35 minutes and notebook 04 about 15 minutes.
\end{itemize}

\section{Team Contributions}
\begin{table}[H]
\centering\small
\caption{Team T068 (Section B): members and primary responsibilities in Phase 2. Every member can explain the full
pipeline.}
\begin{tabularx}{\linewidth}{@{}l l X@{}}
\toprule
\textbf{Name} & \textbf{Enrollment No.} & \textbf{Primary responsibility (Phase 2)}\\
\midrule
Sharma Piyush & 2401010437 & DL data handling: legitimate-only training, label-free holdout, contamination study\\
Saransh Singh Tomar & 2401010418 & Input preprocessing for the autoencoder and the scaling ablation\\
Satwik Mani Tripathi & 2401020056 & Hybrid second stage: latent GMMs, augmented GMM and the evidence analysis\\
Gourav Tanwar & 2401010173 & Autoencoder and denoising autoencoder: architecture, ablations and model selection\\
Alok Kumar & 2401010058 & Evaluation, significance tests, label efficiency, report and presentation\\
\bottomrule
\end{tabularx}
\end{table}

\end{document}
